\chapter{Ciała różniczkowe i twierdzenie Liouville'a}

\section{Pochodna, stałe i lokalny sens funkcji}

Pytanie o wzór przez pierwiastniki dotyczy równań algebraicznych. Pytanie o funkcję pierwotną dotyczy równania $DF=f$. Do drugiego problemu potrzebujemy ciała wyposażonego także w operację różniczkowania. Kluczowym narzędziem będzie twierdzenie Liouville'a, a nie samo kryterium rozwiązalności skończonej grupy Galois.

\begin{definicja}[Ciało różniczkowe]
Ciało różniczkowe $(K,D)$ jest ciałem charakterystyki zero z derywacją $D:K\to K$, to znaczy operacją spełniającą
\[
D(a+b)=Da+Db,\qquad D(ab)=a\,Db+b\,Da.
\]
Ciało stałych to $C_K=\{a\in K:Da=0\}$. Rozszerzenie różniczkowe zachowuje derywację na ciele bazowym. Jest \emph{rozszerzeniem bez nowych stałych}, jeśli ciało stałych się nie powiększa.
\end{definicja}

Z reguły iloczynu wynikają $D1=0$ i $D(a^{-1})=-Da/a^2$ dla $a\ne0$, więc otrzymujemy zwykłą regułę ilorazu. Podstawowy przykład to $\C(x)$ z $D=\mathrm{d}/\mathrm{d}x$ i stałymi $\C$. W całej tej części zakładamy algebraiczną domkniętość ciała stałych $C$; w przykładach będzie to właśnie $\C$.

W interpretacji analitycznej pracujemy lokalnie, na małym otoczeniu punktu, w którym wybrane gałęzie pierwiastków i logarytmów są określone. Elementami ciała mogą być zarodki funkcji meromorficznych: dwa wyrażenia reprezentują ten sam zarodek, jeśli zgadzają się w pewnym otoczeniu punktu. Zarodek o zerowej pochodnej jest stały. Ciała takich funkcji, zawierające wszystkie stałe zespolone, nie dodają więc nowych stałych. Nie żądamy globalnej jednowartościowości logarytmu.

\begin{definicja}[Rozszerzenie elementarne]
Rozszerzenie $E/K$ jest elementarne, gdy istnieje skończona wieża $K=K_0\subseteq\cdots\subseteq K_r=E$, w której $K_i=K_{i-1}(\theta_i)$ i zachodzi jeden z warunków:
\begin{enumerate}
\item $\theta_i$ jest algebraiczny nad $K_{i-1}$;
\item $D\theta_i=Du/u$ dla pewnego $u\in K_{i-1}^\times$ (krok logarytmiczny);
\item $\theta_i\ne0$ i $D\theta_i/\theta_i=Du$ dla pewnego $u\in K_{i-1}$ (krok wykładniczy).
\end{enumerate}
W twierdzeniu o całkowaniu wymagamy dodatkowo $C_E=C_K$. Element $f\in K$ ma elementarną funkcję pierwotną nad $K$, jeśli w takim rozszerzeniu istnieje $F$ spełniające $DF=f$.
\end{definicja}

To definicja obejmująca działania arytmetyczne, funkcje algebraiczne, wykładnicze i logarytmy, zagnieżdżane skończenie wiele razy. Funkcje trygonometryczne też się w niej mieszczą, na przykład $\sin x=(e^{ix}-e^{-ix})/(2i)$. Także $\arctan x=(\log(1+ix)-\log(1-ix))/(2i)$, po lokalnym wyborze gałęzi i stałej. Nie wolno natomiast zastępować drugiego warunku dowolnym $D\theta\in K_{i-1}$: oznaczałoby to dopuszczenie dowolnych całek już w definicji, a więc zmianę badanego problemu.

\section{Krok algebraiczny: pochodna, ślad i norma}

\begin{lemat}
Derywacja ciała charakterystyki zero przedłuża się jednoznacznie do każdego jego rozszerzenia algebraicznego. Jeśli ciało stałych $C$ jest algebraicznie domknięte, takie rozszerzenie nie dodaje stałych.
\end{lemat}
\begin{proof}
Niech $\alpha$ ma unormowany wielomian minimalny $p(T)\in K[T]$. Przez $p^D$ oznaczamy wielomian otrzymany przez zróżniczkowanie współczynników, a przez $p_T$ zwykłą pochodną względem $T$. Z relacji $p(\alpha)=0$ musi wynikać
\[
D\alpha=-\frac{p^D(\alpha)}{p_T(\alpha)}.
\]
Mianownik jest niezerowy dzięki rozdzielczości. Wzór $D(q(\alpha))=q^D(\alpha)+q_T(\alpha)D\alpha$ jest dobrze określony: pochodna wielokrotności $p$ znika po podstawieniu $\alpha$ właśnie dzięki powyższej równości. Reguły sumy i iloczynu sprawdzamy na wielomianach, a reguła ilorazu przedłuża derywację na ciało. To dowodzi istnienia i jednoznaczności na rozszerzeniu prostym. Indukcja daje przedłużenie na każde rozszerzenie skończone. Każdy element dowolnego rozszerzenia algebraicznego należy do pewnego skończonego podrozszerzenia. Na dwóch takich podrozszerzeniach pochodne obliczone z wielomianu minimalnego są zgodne na części wspólnej, ponieważ przedłużenie derywacji do ciała algebraicznego, jeśli istnieje, jest jednoznaczne. Definiuje to derywację na sumie skierowanej wszystkich skończonych podrozszerzeń, czyli na całym rozszerzeniu algebraicznym.

Jeżeli algebraiczny element $\alpha$ spełnia $D\alpha=0$, to $p^D(\alpha)=0$. Wielomian $p$ jest unormowany, więc $\deg p^D<\deg p$, o ile $p^D\ne0$. Minimalność wymusza $p^D=0$. Wszystkie współczynniki $p$ należą do $C$, zatem $\alpha$ jest algebraiczny nad $C$. Domkniętość algebraiczna $C$ daje $\alpha\in C$.
\end{proof}

Zanurzenia algebraiczne nad ciałem różniczkowym komutują z $D$: oba sposoby obliczenia pochodnej obrazu generatora dają ten sam wzór z wielomianu minimalnego. Dla skończonego $L/K$, różniczkując sumę i iloczyn po wszystkich zanurzeniach, otrzymujemy więc
\begin{equation}\label{eq:trace-norm}
D\Tr_{L/K}(v)=\Tr_{L/K}(Dv),\qquad
\frac{D\Nm_{L/K}(u)}{\Nm_{L/K}(u)}
=\Tr_{L/K}\!\left(\frac{Du}{u}\right),\quad u\ne0.
\end{equation}
W drugiej równości użyliśmy reguły pochodnej iloczynu podzielonej przez sam iloczyn. Normalność $L/K$ nie jest wymagana.

\section{Część wielomianowa funkcji wymiernej}

Niech $t$ będzie przestępny nad ciałem różniczkowym $F$. Dzielenie wielomianów daje jednoznaczny rozkład $R=P+Q$ dla $R\in F(t)$, gdzie $P\in F[t]$, a $Q$ jest właściwą funkcją wymierną, czyli ma licznik stopnia mniejszego od mianownika. Oznaczamy $P=\operatorname{pp}_t(R)$.

Równoważnie rozwijamy formalnie względem $t^{-1}$:
\[
R=\sum_{j\le m}r_jt^j,\qquad r_j\in F,\quad
\operatorname{pp}_t(R)=\sum_{j\ge0}r_jt^j.
\]
Rozwinięcie istnieje, ponieważ po wyłączeniu potęgi $t$ z mianownika możemy odwrócić szereg o niezerowym wyrazie stałym. Jest jednoznaczne. To rachunek formalny w zmiennej $t$, nie granica $x\to\infty$ po podstawieniu $t=e^x$.

\begin{lemat}[Części wielomianowe pochodnych]
Jeśli $Dt=a\in F$ albo $Dt=ht$ z $h\in F$, to
\[
\operatorname{pp}_t(DR)=D\operatorname{pp}_t(R).
\]
Dla $u=b\,p(t)/q(t)\ne0$, gdzie $b\in F^\times$, a $p,q\in F[t]$ są unormowanymi wielomianami, zachodzi
\[
\operatorname{pp}_t\left(\frac{Du}{u}\right)=
\begin{cases}
Db/b, & Dt=a,\\
Db/b+(\deg p-\deg q)h, & Dt=ht.
\end{cases}
\]
W szczególności te pochodne logarytmiczne nie mają dodatnich potęg $t$ w części wielomianowej.
\end{lemat}
\begin{proof}
Dla pojedynczego wyrazu szeregu mamy odpowiednio
\[
D(r_jt^j)=(Dr_j)t^j+j a r_jt^{j-1},\qquad
D(r_jt^j)=(Dr_j+jhr_j)t^j.
\]
Ujemne potęgi pozostają ujemne, a pochodna wielomianu pozostaje wielomianem. To daje pierwszą równość.

Dla unormowanego $p$ stopnia $d$ i $Dt=a$ zachodzi $\deg Dp<d$. Dlatego $Dp/p$ jest właściwe. Dla $Dt=ht$ współczynnik wiodący $Dp$ wynosi $dh$, więc $Dp/p-dh$ jest właściwe. Stosujemy te obserwacje do tożsamości $Du/u=Db/b+Dp/p-Dq/q$.
\end{proof}

\section{Kiedy dołączenie funkcji pierwotnej dodaje stałe?}

\begin{lemat}
Niech $C_F=C$, a $t$ będzie przestępny nad $F$ z $Dt=a\in F$. Wtedy $C_{F(t)}=C$ dokładnie wtedy, gdy w $F$ nie ma elementu $b$ spełniającego $Db=a$.
\end{lemat}
\begin{proof}
Jeśli $Db=a$, to $t-b$ jest nową stałą, nienależącą do $F$. Odwrotnie, niech $R\in F(t)\setminus C$ i $DR=0$. Rozwińmy $R$ formalnie w nieskończoności. Jeżeli najwyższy wykładnik wynosi $0$, to jego współczynnik ma pochodną zero i należy do $C$; odejmujemy go. Pozostaje niezerowy szereg o najwyższym wykładniku $m\ne0$, ponieważ w przeciwnym razie $R$ od początku należałoby do $C$. Możliwy jest również $m<0$.

Z dwóch najwyższych współczynników równania $DR=0$ otrzymujemy
\[
Dr_m=0,\qquad Dr_{m-1}+m a r_m=0,\qquad r_m\ne0.
\]
W charakterystyce zero liczba $m$ jest niezerowa w $F$. Ponieważ $r_m\in C^\times$, element $b=-r_{m-1}/(mr_m)\in F$ spełnia $Db=a$. To dowodzi koniecznego i dostatecznego warunku.
\end{proof}

\section{Twierdzenie Liouville'a o całkowaniu elementarnym}

\begin{twierdzenie}[Liouville'a]
Niech $(K,D)$ ma charakterystykę zero i algebraicznie domknięte ciało stałych $C$. Element $f\in K$ ma funkcję pierwotną w elementarnym rozszerzeniu bez nowych stałych wtedy i tylko wtedy, gdy istnieją $v\in K$, $u_1,\dots,u_r\in K^\times$, $c_1,\dots,c_r\in C$ takie, że
\begin{equation}\label{eq:liouville}
f=Dv+\sum_{j=1}^r c_j\frac{Du_j}{u_j}.
\end{equation}
Dopuszczamy $r=0$. Wszystkie $v,u_j$ można wybrać w ciele bazowym $K$, a nie dopiero w jego rozszerzeniu algebraicznym.
\end{twierdzenie}
\begin{proof}
\emph{Konieczność: schodzenie po wieży.} W końcowym ciele wieży elementarnej mamy $f=DV$, czyli przedstawienie \eqref{eq:liouville} bez składników logarytmicznych. Pokażemy, że przedstawienie z $v,u_j$ w pojedynczym rozszerzeniu $L/F$ można zastąpić takim przedstawieniem z elementami $F$, jeśli $f\in F$. Wszystkie ciała pośrednie mają stałe $C$, ponieważ zawierają $C$ i są podciałami końcowego rozszerzenia bez nowych stałych.

\emph{1. Krok algebraiczny.} Jeśli $[L:F]=d$, bierzemy ślad przedstawienia i dzielimy przez $d$. Równości \eqref{eq:trace-norm} dają
\[
f=D\left(\frac{\Tr_{L/F}(v)}{d}\right)
+\sum_j\frac{c_j}{d}\,
\frac{D\Nm_{L/F}(u_j)}{\Nm_{L/F}(u_j)}.
\]
Ślad i normy należą do $F$, normy są niezerowe, a $c_j/d\in C$. Otrzymujemy przedstawienie na niższym piętrze.

\emph{2. Przestępny krok logarytmiczny.} Niech $L=F(t)$, $Dt=a=Db/b$, $b\in F^\times$. Zastosujmy część wielomianową względem $t$ do przedstawienia $f$. Jeśli $P=\operatorname{pp}_t(v)$, lemat 8.4 daje
\[
f=DP+\sum_j c_j\frac{Db_j}{b_j},\qquad b_j\in F^\times.
\]
W szczególności $DP\in F$. Gdy $P=p_mt^m+\cdots+p_0$ i $m\ge1$, współczynnik przy $t^m$ daje $Dp_m=0$, więc $p_m\in C^\times$. Jeśli $m\ge2$, musi zniknąć także współczynnik przy $t^{m-1}$:
\[
Dp_{m-1}+m p_ma=0.
\]
Stąd $a=Dw$ dla $w=-p_{m-1}/(mp_m)\in F$. Wówczas $D(t-w)=0$, więc brak nowych stałych dawałby $t-w\in C\subseteq F$, sprzecznie z przestępnością $t$. Pozostaje $P=p_0+ct$ z $c\in C$ (także $c=0$). Wobec tego
\[
f=Dp_0+c\frac{Db}{b}+\sum_j c_j\frac{Db_j}{b_j},
\]
co jest wymaganą postacią w $F$.

\emph{3. Przestępny krok wykładniczy.} Niech $L=F(t)$, $Dt=ht$ z $h=Da$, $a\in F$. Część wielomianowa daje teraz
\[
f=DP+\sum_j c_j\frac{Db_j}{b_j}+Ah,\qquad A=\sum_j c_jm_j\in C,
\]
gdzie $m_j\in\Z$ są różnicami stopni liczników i mianowników. Dla $P=\sum_{j\ge0}p_jt^j$ współczynniki dodatnich potęg spełniają $Dp_j+jhp_j=0$. Niezerowe $p_j$ dla $j>0$ dawałoby $D(p_jt^j)=0$, więc $p_jt^j\in C^\times$. Oznaczałoby to $t^j\in F$, wbrew przestępności $t$. Zatem $P=p_0$ i
\[
f=D(p_0+Aa)+\sum_j c_j\frac{Db_j}{b_j}.
\]
Tu wykorzystaliśmy dokładnie warunek $h=Da$ z definicji kroku wykładniczego.

Jeśli generator logarytmiczny lub wykładniczy jest algebraiczny, stosujemy krok 1 zamiast kroków 2--3. Przechodząc skończenie wiele razy w dół wieży, otrzymujemy \eqref{eq:liouville} w $K$.

\emph{Dostateczność: kontrolowane dołączanie logarytmów.} Mając \eqref{eq:liouville}, chcemy otrzymać $v+\sum_j c_j\log u_j$. Aby nie wprowadzić nowych stałych formalnym dołączeniem zbędnego generatora, postępujemy kolejno. W aktualnym ciele szukamy funkcji pierwotnej $b_j$ elementu $Du_j/u_j$. Jeśli istnieje, korzystamy z niej bez powiększania ciała. Jeśli nie istnieje, dołączamy przestępny $z_j$ z relacją $Dz_j=Du_j/u_j$. Lemat 8.5 zapewnia, że nie powstają nowe stałe. Każdy dołączany krok jest logarytmiczny. Po skończenie wielu krokach element $v+\sum_j c_jb_j$ (gdzie $b_j$ oznacza wybraną lub dołączoną funkcję pierwotną) ma pochodną $f$, co kończy dowód.
\end{proof}

\begin{uwaga}
Twierdzenie ogranicza postać jakiejś funkcji pierwotnej, o ile istnieje elementarna funkcja pierwotna; nie wymaga, by pierwotnie podany wzór od razu był zapisany w tej postaci. Przy algebraicznie domkniętych stałych przedstawienie w skończonym rozszerzeniu algebraicznym zawsze można sprowadzić do bazy śladem i normą. Nie jest to dodatkowe założenie dotyczące wybranych przykładów.
\end{uwaga}

\section{Sprawdzenie założeń dla funkcji wykładniczych}

\begin{lemat}
Jeżeli $p\in\C[x]$ jest niestały, to $t=e^{p(x)}$ jest przestępny nad $\C(x)$, a ciało stałych $\C(x,t)$ jest równe $\C$.
\end{lemat}
\begin{proof}
Dla dowolnego $q\in\C(x)^\times$ rozkład na czynniki liniowe daje
\[
q=c\prod_j(x-a_j)^{m_j},\qquad
\frac{q'}{q}=\sum_j\frac{m_j}{x-a_j},\qquad m_j\in\Z.
\]
Suma jest skończona; wykładniki ujemne oznaczają bieguny. Pochodna logarytmiczna $q'/q$ jest więc właściwą funkcją wymierną i nie może być niezerowym wielomianem.

Gdyby $t$ było algebraiczne, dla $L=\C(x,t)$ stopnia $d$ norma $q=\Nm_{L/\C(x)}(t)\ne0$ spełniałaby
\[
\frac{q'}{q}=\Tr_{L/\C(x)}(Dt/t)=d\,p'(x).
\]
Prawa strona jest niezerowym wielomianem, co daje sprzeczność.

Pozostaje sprawdzić stałe; można to zrobić także czysto algebraicznie. Gdyby $R\in\C(x,t)\setminus\C$ spełniało $DR=0$, w formalnym rozwinięciu względem $t^{-1}$ wybieramy najwyższą potęgę $t^m$. Jeśli $m=0$, współczynnik jest stały; odejmujemy go i otrzymujemy niezerowy szereg z $m\ne0$. Współczynnik $r_m\ne0$ musi spełniać
\[
\frac{r_m'}{r_m}=-m p'(x).
\]
Znów pochodna logarytmiczna funkcji wymiernej musiałaby być niezerowym wielomianem. To niemożliwe, więc nowych stałych nie ma.
\end{proof}

\section{Kryterium hiperwykładnicze z dowodem}

\begin{definicja}[Element hiperwykładniczy]
W tym skrypcie elementem hiperwykładniczym nad $\C(x)$ nazywamy element $t$ przestępny nad $\C(x)$, spełniający $Dt=h(x)t$ dla $h\in\C(x)$, przy założeniu $C_{\C(x,t)}=\C$. Poprzedni lemat sprawdza te założenia dla wszystkich $e^{p(x)}$ z niestałym wielomianem $p$.
\end{definicja}

\begin{twierdzenie}[Kryterium hiperwykładnicze]
Niech $t$ będzie jak w definicji, a $s\in\C(x)$. Element $s(x)t$ ma elementarną funkcję pierwotną nad $\C(x,t)$ wtedy i tylko wtedy, gdy istnieje $r\in\C(x)$ spełniające
\begin{equation}\label{eq:hyperexp}
r'+hr=s.
\end{equation}
Jeżeli rozwiązanie istnieje, funkcją pierwotną jest $rt$.
\end{twierdzenie}
\begin{proof}
Dostateczność wynika z $D(rt)=(r'+hr)t=st$. Przy elementarnej funkcji pierwotnej stosujemy twierdzenie Liouville'a z ciałem bazowym $\C(x,t)$. Otrzymujemy
\[
st=Dv+\sum_j c_j\frac{Du_j}{u_j},\qquad v,u_j\in\C(x,t).
\]
Pochodne logarytmiczne nie mają dodatnich potęg $t$ w części wielomianowej, na mocy lematu 8.4. Niech $r$ będzie współczynnikiem przy $t$ w $\operatorname{pp}_t(v)$, z $r=0$, jeśli takiego wyrazu nie ma. Porównanie współczynnika przy $t$ daje dokładnie $s=r'+hr$. Przestępność $t$ uzasadnia jednoznaczność tego porównania.
\end{proof}

\begin{wniosek}[Rozdzielenie potęg w sumie Laurenta]
Jeżeli skończona suma $\sum_k s_k(x)t^k$, $s_k\in\C(x)$, ma elementarną funkcję pierwotną, to dla każdego $k\ne0$ istnieje $r_k\in\C(x)$ spełniające
\[
r_k'+khr_k=s_k.
\]
\end{wniosek}
\begin{proof}
Dla $k>0$ stosujemy poprzednie porównanie do współczynnika przy $t^k$. Dla $k<0$ zmieniamy generator na $w=t^{-1}$. Mamy to samo ciało $\C(x,w)=\C(x,t)$ i $Dw=-hw$, więc wystarczy ponownie porównać dodatnie potęgi nowego generatora.
\end{proof}

Nie można po prostu porównywać ujemnych potęg w tym samym rozwinięciu w nieskończoności: na przykład
\[
\frac{D(t-1)}{t-1}=\frac{ht}{t-1}
=h+\frac{h}{t-1}=h+ht^{-1}+ht^{-2}+\cdots.
\]
Pochodna logarytmiczna nie ma tu dodatnich potęg, ale ma ujemne. Dla $k=0$ wniosek także nie twierdzi istnienia wymiernej funkcji pierwotnej: funkcja $1/x$ ma funkcję pierwotną $\log x$, a nie funkcję wymierną. W ten sposób redukcja nie zakłada żadnej nieprawdziwej niezależności $t$ i $t^{-1}$.


